\section{Use a conservative capacity before reading pixels}
For valid dimensions with \(n=wh\), define the encoder bound
\[
B(n)=22+5n.
\]
The completed encoder requires a destination capacity of at least this bound before
reading input pixels or writing caller data. Return a capacity error for a smaller
capacity even when a particular image could compress into it. This is a deliberate
interface choice that makes failure independent of a preliminary scan of the pixels.
The successful byte count remains the actual encoded length, not the bound.

For the eight-by-eight sprite, the bound is 342 bytes. The worked compressed fixture
uses sixty-one bytes, while its expanded RGBA words occupy 256 bytes. Its original
monochrome row representation still occupies only eight bytes. QOI carries a general
image header and supports many colors; it is not expected to beat a representation
specialized to one-bit pixels in every small example.

For a sixty-four-by-sixty-four image, the bound is 20502 bytes. Validate width and height
before multiplication, then perform the arithmetic in registers that retain all the
required bits. A bound helper can accept packed dimensions directly and return this
size without touching caller memory. The encoder additionally validates the descriptor's
channels and colorspace fields before using its pixels.

Do not fill unused destination capacity with zeros. Only the actual header, chunks
and marker form the returned output interval. A trailing terminator or an extra marker
byte changes the file and violates the exact-extent contract.
